% TOP_PROBLEM: {"id":30007554,"problem_number":"LOCAL-30007554","title":"Welfare from inconsistent choices","category_id":21,"category":{"id":21,"name":"economics","display_name":"Economics","description":"Open economic theory statements, published research agendas, and proposed scoped specifications; question provenance and historical priority are qualified per record.","slug":"economics","order_index":21,"created_at":"2026-10-08T00:00:00Z"},"status":"research_question","external_url":null,"legacy_ids":[],"published":false,"statement":"Determine the sharp welfare comparisons shared by all maintained behavioral models consistent with choice and process data, and which experiments can reduce the remaining ambiguity.","economics":{"original_id":43,"field":"Economic theory, equilibrium and social choice","kind":"identification","formalization_basis":"adapted_research_question","source_status":"research_agenda","priority_status":"earliest_found","priority_note":"This is the earliest source in which the relevant agenda was verified during this audit; absolute priority has not been established.","earliest_verified_reference":{"authors":["Geoffroy de Clippel","Kareen Rozen"],"title":"Bounded Rationality in Choice Theory: A Survey","year":2023,"url":"https://geoffroydeclippel.net/Working%20Papers%20PDFs/JEL.pdf","doi":"10.1257/jel.20231592","locator":"§6.4 Searching for common ground; §7 Looking Forward, PDF pp.51–54","evidence_summary":"The revised working paper identifies welfare comparisons, additional process data, and computational or procedural constraints as continuing research directions."},"current_reference":{"authors":["Geoffroy de Clippel","Kareen Rozen"],"title":"Bounded Rationality in Choice Theory: A Survey","year":2023,"url":"https://geoffroydeclippel.net/Working%20Papers%20PDFs/JEL.pdf","doi":"10.1257/jel.20231592","locator":"§6.4 Searching for common ground; §7 Looking Forward, PDF pp.51–54","evidence_summary":"The revised working paper identifies welfare comparisons, additional process data, and computational or procedural constraints as continuing research directions."},"formalization_note":"The robust intersection and statistical target are specified here; the survey motivates model-dependent welfare and additional evidence rather than stating this exact theorem."},"statusQualification":"Published question or research agenda verified at the cited locator. The mathematical specification is adapted for this catalogue and includes additional assumptions; source publication does not establish first-ever priority or certify this exact model remains unsolved. The robust intersection and statistical target are specified here; the survey motivates model-dependent welfare and additional evidence rather than stating this exact theorem.","statusReviewedAt":"2026-10-08"}
% FIRST_POSTED: 2026-10-08
% ENTRY_KIND: statement_only
% !TeX program = lualatex
\documentclass[11pt]{article}
\usepackage{fontspec}
\usepackage[a4paper,margin=25mm]{geometry}
\usepackage{amsmath,amssymb,mathtools}
\usepackage{xurl}
\usepackage[hidelinks]{hyperref}
\newcommand{\sref}[2]{\hyperlink{source:#1:#2}{[S#2]}}
\setlength{\emergencystretch}{3em}
\begin{document}
% problemId:30007554
\section*{Welfare from inconsistent choices}\label{30007554}
\subsection{Definitions and mathematical statement}
Let \(X\) be a finite set of economic alternatives, \(B\subseteq X\) a feasible menu, and \(f\in F\) an observed frame. For each menu and frame, observe choice probabilities \(P(x\mid B,f)\) and optional process measurements \(L\), such as response time or attention. Fix a publicly specified class \(\mathcal M\) of behavioral models. A model \(m\in\mathcal M\) consists of a latent welfare ordering \(\succeq_m\), a choice procedure, and a distribution of \(L\). State explicitly the normative assumption connecting that ordering to welfare; inconsistent choices alone cannot supply it. Let \(\mathcal M(P,L)\) contain all models fitting the population observables. Define robust welfare comparison by
\[
 x\succeq^R y\quad\Longleftrightarrow\quad x\succeq_m y\text{ for every }m\in\mathcal M(P,L).
\]
Determine the sharp identified relation \(\succeq^R\), which additional experimentally manipulated frames or process measurements enlarge it, and which comparisons remain intrinsically ambiguous within the maintained model class. Construct uncertainty regions accounting for sampling error and evaluate policy rankings by their robust welfare implications. A successful solution must report how conclusions change when competing attention or preference-construction models are admitted. This formalizes a research agenda for combining welfare assumptions with richer behavioral evidence, rather than asserting a unique welfare ordering without normative restrictions.

\par\textbf{Formulation and scope.} The robust intersection and statistical target are specified here; the survey motivates model-dependent welfare and additional evidence rather than stating this exact theorem.

\par\textbf{Open-question provenance.} An explicit question or research agenda was verified at the source locator. This does not by itself certify that every later version remains unresolved \sref{30007554}{1}.

\par\textbf{Historical priority.} This is the earliest source in which the relevant agenda was verified during this audit; absolute priority has not been established.

\subsection{Short English statement}
Determine the sharp welfare comparisons shared by all maintained behavioral models consistent with choice and process data, and which experiments can reduce the remaining ambiguity.

\subsection{Sources}
\begin{itemize}\raggedright
\item[\textnormal{[S1]}]\hypertarget{source:30007554:1}{} Geoffroy de Clippel, Kareen Rozen. 2023. Bounded Rationality in Choice Theory: A Survey. §6.4 Searching for common ground; §7 Looking Forward, PDF pp.51–54.
\url{https://geoffroydeclippel.net/Working%20Papers%20PDFs/JEL.pdf}.
DOI: 10.1257/jel.20231592.
{\small\textit{Earliest verified question/agenda reference; historical priority qualified below. The revised working paper identifies welfare comparisons, additional process data, and computational or procedural constraints as continuing research directions.}}
\end{itemize}
\end{document}
